\chapter{Auctions}\label{ch:m1:auctions}

At ten minutes to four in New York the order book of a large stock looks as
it has all day: a few hundred shares at each price, a cent apart. Beside it,
invisible on the ordinary screen, orders to buy and sell several million
shares have accumulated with a single instruction: \emph{at the \omterm{def:m1:auctions:moc}{closing
price}, whatever it is}. At four o'clock they will all trade, at one price,
in one instant, in a quantity the continuous market could not have absorbed
in an hour. That price will value every fund holding the stock, settle the
options that expire today and become the reference for tomorrow's limits. It
is computed by an algorithm of about thirty lines. This chapter is that
algorithm, the information published while it waits, and the reasons one
print a day has become the most important one.

\section{The call auction}

\begin{definition}[Call auction and uncrossing price]\label{def:m1:auctions:call}
In a \emph{call auction}\index{call auction} orders are collected for a period
without being executed; at the end of it a single
\emph{uncrossing price}\index{uncrossing price} is computed, and every order
that can trade at that price does, at that price. Opening and closing
auctions, \omterm{def:m1:auctions:volint}{volatility interruptions} and initial public offerings are call
auctions; so, in miniature, are the \omterm{def:m1:european-equity-market-structure:periodic}{periodic auctions} of
\cref{def:m1:european-equity-market-structure:periodic}.
\end{definition}

For a candidate price $p$ write $D(p)$ for the total quantity of buy orders
willing to pay $p$ or more (market orders included) and $S(p)$ for the total
quantity of sell orders willing to accept $p$ or less. $D$ is decreasing, $S$
increasing; the quantity that can trade at $p$ is $\min(D(p), S(p))$.

\begin{method}[The uncrossing rules]\label{met:m1:auctions:rules}
Among the limit prices present in the book and the reference price:
\begin{enumerate}
\item \textbf{Volume.} Keep the prices that maximise the executable volume
  $\min(D(p),S(p))$.
\item \textbf{Surplus.} Among those, keep the prices that minimise the
  unexecuted quantity $|D(p) - S(p)|$.
\item \textbf{Market pressure.} If at every remaining price the surplus is on
  the \omterm{def:m1:the-buy-side:buy-side}{buy side}, take the highest; if on the \omterm{def:m1:the-sell-side:sell-side}{sell side}, the lowest.
\item \textbf{Reference price.} Otherwise take the price closest to the
  reference price (usually the last traded price).
\end{enumerate}
Then allocate: market orders first, then limit orders priced better than the
\omterm{def:m1:auctions:call}{uncrossing price}, then orders at the \omterm{def:m1:auctions:call}{uncrossing price} in time priority, until
the executable volume is exhausted on each side.
\end{method}

The first two rules are common to the European exchanges' published market
models: the auction price is the one ``with the most executable volume and
the lowest surplus''. The tie-breaking rules vary in detail from venue to
venue, and an implementation must follow the rulebook of the venue it
simulates; the logic above is the usual one.

\begin{example}[Eight orders]\label{ex:m1:auctions:eight}
Buyers: 300 at market; 500 at 10.03; 400 at 10.01; 600 at 10.00. Sellers: 200
at market; 400 at 9.99; 500 at 10.01; 700 at 10.02. The reference price is
10.00.
\begin{center}\small
\begin{tabular}{@{}lrrrr@{}}
\hline
Price & $D(p)$ & $S(p)$ & Executable & Surplus \\
\hline
9.99 & 1\,800 & 600 & 600 & $+1\,200$ \\
10.00 & 1\,800 & 600 & 600 & $+1\,200$ \\
10.01 & 1\,200 & 1\,100 & \textbf{1\,100} & $+100$ \\
10.02 & 800 & 1\,800 & 800 & $-1\,000$ \\
10.03 & 800 & 1\,800 & 800 & $-1\,000$ \\
\hline
\end{tabular}
\end{center}
Rule 1 alone selects 10.01: 1\,100 shares trade. On the \omterm{def:m1:the-buy-side:buy-side}{buy side} the market
order (300) and the order at 10.03 (500) are filled in full and the order at
10.01 receives 300 of its 400; on the \omterm{def:m1:the-sell-side:sell-side}{sell side} everything up to 10.01 is
filled. A buy surplus of 100 shares remains at 10.01, and that is what the
exchange publishes as the closing imbalance.
\end{example}

\begin{omfigure}
\begin{tikzpicture}
\begin{axis}[width=10.5cm,height=5.2cm,
  xlabel={price},ylabel={quantity (shares)},
  tick label style={font=\small},label style={font=\small},
  xmin=9.965,xmax=10.055,ymin=0,ymax=2000,grid=major,grid style={gray!25},
  xtick={9.97,9.98,9.99,10.00,10.01,10.02,10.03,10.04,10.05},
  xticklabel style={/pgf/number format/fixed,/pgf/number format/precision=2,/pgf/number format/fixed zerofill},
  yticklabel style={/pgf/number format/1000 sep={\,}},
  legend columns=2,legend style={font=\footnotesize,at={(0.5,-0.3)},anchor=north,draw=none,fill=none,/tikz/every even column/.append style={column sep=8pt}},
  table/col sep=comma,x filter/.code={\pgfmathparse{\pgfmathresult/100}}]
\addplot[omDef,very thick,const plot mark mid,mark=*,mark size=1.3pt] table[x=price,y=demand]{figdata/markets-1/13-auctions/curves.csv};
\addplot[omThm,very thick,const plot mark mid,mark=*,mark size=1.3pt] table[x=price,y=supply]{figdata/markets-1/13-auctions/curves.csv};
\addplot[only marks,mark=o,mark size=4pt,black,thick] coordinates {(1001,1100)};
\legend{demand $D(p)$,supply $S(p)$}
\end{axis}
\end{tikzpicture}

\omcaption{Demand and supply of \cref{ex:m1:auctions:eight}. The curves are
step functions of the price; the auction trades the largest quantity on which
they can agree, 1\,100 shares at 10.01 (circled: the lower of the two curves is
highest there). Data: the chapter's
example book.}\label{fig:m1:auctions:curves}
\end{omfigure}

\begin{proposition}[Why a single price]\label{prop:m1:auctions:welfare}
Among all sets of trades between the orders in the book that respect each
order's limit, the uncrossing at a price satisfying rule 1 maximises the
traded volume, and every participant trades at a price at least as good as
its limit with no participant able to complain that a worse-priced order
traded in its place.
\end{proposition}
\begin{proof}
Any feasible set of trades at arbitrary prices pairs a buyer with limit $b$
and a seller with limit $s \le b$. Order buyers by decreasing limit and
sellers by increasing limit: the $k$-th pair is feasible iff the $k$-th best
bid is at least the $k$-th best offer, and the largest such $k$ (in shares)
is $\max_p \min(D(p),S(p))$, attained by pairing in that order. Trading all
of them at one price $p^\star$ between the marginal bid and the marginal
offer respects every limit, and any excluded order has a limit no better than
an included one.
\end{proof}

\section{Opening, closing and everything else}

\begin{definition}[Closing price, market-on-close and limit-on-close orders]\label{def:m1:auctions:moc}
The \emph{closing price}\index{closing price} of a security is the price of
its closing auction on its listing exchange. A
\emph{market-on-close order}\index{market-on-close order} (MOC) executes in
that auction at whatever price it sets; a
\emph{limit-on-close order}\index{limit-on-close order} (LOC) executes in it
only at its limit or better.
\end{definition}

\begin{definition}[Order imbalance and indicative price]\label{def:m1:auctions:imbalance}
During the call period the exchange publishes the
\emph{indicative price}\index{indicative price}, the price at which the
auction would uncross if it ended now, the quantity \emph{paired} at that
price, and the \emph{order imbalance}\index{order imbalance}: the side and
size of the quantity that could not be paired.
\end{definition}

\begin{dated}{2026-09}{The last ten minutes in New York}\label{dat:m1:auctions:us}
On the New York Stock Exchange, MOC and LOC orders may be entered until
15:50; at that time the exchange publishes a regulatory imbalance for every
symbol whose imbalance is at least 500 \omterm{def:m1:us-equity-market-structure:lots}{round lots}, after which closing orders
are accepted only on the opposite side, and imbalance information is
republished every second. On Nasdaq, imbalance information is disseminated
from 15:50, every five seconds and then every second from 15:55; MOC orders
are accepted until 15:55 and LOC orders until 15:58. Both auctions run at
16:00.
\end{dated}

\begin{omfigure}
\begin{tikzpicture}[font=\small]
\draw[-{Stealth[length=2.2mm]},thick,omExo] (-0.4,0) -- (15.2,0);
\foreach \x/\t in {1.0/15:50,6.0/15:55,10.0/15:58,14.0/16:00}{
  \draw[omExo,thick] (\x,-0.12) -- (\x,0.12);
  \node[font=\small\bfseries,anchor=north] at (\x,-0.2) {\t};
}
\node[omDef,font=\footnotesize,align=center,anchor=south] at (1.0,0.25) {imbalance published;\\NYSE closes MOC and\\LOC entry};
\node[omDef,font=\footnotesize,align=center,anchor=south] at (6.0,0.25) {Nasdaq closes\\MOC entry};
\node[omDef,font=\footnotesize,align=center,anchor=south] at (10.0,0.25) {Nasdaq closes\\LOC entry};
\node[omThm,font=\footnotesize,align=center,anchor=south] at (14.0,0.25) {both auctions\\uncross};
\draw[omThm,thick,decorate,decoration={brace,mirror,amplitude=5pt}] (1.0,-0.95) -- (14.0,-0.95)
  node[midway,below=6pt,font=\footnotesize,omThm] {imbalance information is published; the continuous market still trades};
\end{tikzpicture}

\omcaption{The closing ten minutes on the two US listing exchanges, from
\cref{dat:m1:auctions:us}. The cut-offs exist so that the published imbalance
means something: after an exchange's cut-off, new closing orders there can
only reduce it.}\label{fig:m1:auctions:timeline}
\end{omfigure}

The \emph{opening} auction solves a different problem: overnight news has
moved the fair price by an unknown amount and the first continuous trades
would be made in the dark. A call period lets everyone see the \omterm{def:m1:auctions:imbalance}{indicative
price} converge before committing. It attracts far less volume than the close,
because nobody is benchmarked to the open.

\begin{definition}[Volatility interruption]\label{def:m1:auctions:volint}
A \emph{volatility interruption}\index{volatility interruption} is an
unscheduled \omterm{def:m1:auctions:call}{call auction} that an exchange starts automatically when the next
trade would occur outside a price corridor around a reference price: the
continuous market stops for a few minutes, orders accumulate, and trading
restarts at an \omterm{def:m1:auctions:call}{uncrossing price}.
\end{definition}

The European exchanges' answer to a sudden move is thus to change mechanism
rather than to forbid prices, as a daily limit does
(\cref{def:m1:asian-and-emerging-equity-markets:limit}), or to halt, as the
American limit-up--limit-down bands do (\cref{ch:m1:stress-case-studies}).

\section{Imbalances are information}

\begin{omfigure}
\begin{tikzpicture}
\begin{axis}[width=10.5cm,height=5.2cm,
  xlabel={buy imbalance added (\% of the balanced volume)},ylabel={move of the close (ticks)},
  tick label style={font=\small},label style={font=\small},
  xmin=0,xmax=100,ymin=0,ymax=4,grid=major,grid style={gray!25},table/col sep=comma]
\addplot[omDef,very thick,mark=*,mark size=1.4pt] table[x=imbalance_pct,y=move_ticks]{figdata/markets-1/13-auctions/impact.csv};
\end{axis}
\end{tikzpicture}

\omcaption{What an unexpected closing buy order does to the price, in
simulated auctions of 400 limit orders. The relation is close to linear
because the supply curve near the \omterm{def:m1:auctions:call}{uncrossing price} is; its slope is the
auction's depth. Data: the tutorial's simulation, 60 books per
point.}\label{fig:m1:auctions:impact}
\end{omfigure}

\begin{proposition}[The price of an imbalance]\label{prop:m1:auctions:depth}
Suppose that near the \omterm{def:m1:auctions:call}{uncrossing price} the book has a constant density: each
tick of price adds $\lambda_S$ shares of supply and removes $\lambda_D$
shares of demand. A market buy order of $q$ shares added to a balanced
auction raises the \omterm{def:m1:auctions:call}{uncrossing price} by approximately $q/(\lambda_D +
\lambda_S)$ ticks.
\end{proposition}
\begin{proof}
At the old price $p_0$, $D = S$. At $p_0 + x$ the new demand is $D - \lambda_D
x + q$ and the supply $S + \lambda_S x$; they are equal at $x = q/(\lambda_D +
\lambda_S)$.
\end{proof}

The denominator is the only thing that protects a \omterm{def:m1:auctions:moc}{closing price} from a large
order, and it is endogenous: it consists of participants who watch the
published imbalance and enter offsetting interest \emph{because} they expect
the imbalance to move the price. Liquidity providers in the close are paid by
the difference between the \omterm{def:m1:auctions:moc}{closing price} and the price a few minutes before
or the next morning; index funds pay it, knowingly, as the price of zero
\omterm{def:m1:the-buy-side:benchmark}{tracking error} (\cref{ch:m1:index-construction-and-rebalancing}). An order
designed to move the close rather than to trade --- ``marking the close'' ---
is market manipulation everywhere, and the closing auction's transparency is
what makes it detectable.

\section{Why the close keeps growing}

Three forces push volume toward the closing auction, and each reinforces the
others. \emph{Benchmarks:} index funds, and every manager measured against an
index, are valued at \omterm{def:m1:auctions:moc}{closing prices}; trading at any other price is a
tracking-error risk taken for no reward (\cref{ch:m1:the-buy-side}).
\emph{Size:} the auction offers the day's deepest liquidity at zero spread,
so large orders that are indifferent to timing prefer it; their presence
makes it deeper still. \emph{Thinning elsewhere:} volume that moves to the
close leaves the continuous session shallower, which raises the cost of
trading intraday and sends more orders to the close. In Europe a quarter of
on-exchange turnover now trades in closing auctions
(\cref{dat:m1:european-equity-market-structure:july}). The listing exchange
owns this auction, and with it one of the last pieces of pricing power in
equity trading.

\section{Tutorial: uncrossing a book}\label{tut:m1:auctions:uncross}

\begin{tutorial}
\textbf{Goal.} Implement the uncrossing rules, reproduce
\cref{ex:m1:auctions:eight}, then watch an \omterm{def:m1:auctions:imbalance}{indicative price} converge.
\textbf{End state:} price 10.01, volume 1\,100, surplus $+100$, and the two
charts of this section.
\begin{tutsteps}
\item \textbf{The four rules}, on integer prices in ticks.
\omcode{code/firm/auction/firm_auction.py}{38}{53}{Uncrossing: volume, surplus, market pressure, reference price.}\label{lst:m1:auctions:uncross}
\item \textbf{Allocation.} Sort each side so that market orders come first,
then better prices, then earlier arrival, and fill until the volume is
exhausted.
\omcode{code/firm/auction/firm_auction.py}{56}{69}{Allocation with price and time priority.}\label{lst:m1:auctions:allocate}
\item \textbf{Run the example}; the acceptance tests also cover an empty
cross, one-sided pressure, a balanced tie broken by the reference price, and
time priority at the \omterm{def:m1:auctions:call}{uncrossing price}.
\item \textbf{A call phase.} Feed 600 random orders in forty batches and
recompute after each: the \omterm{def:m1:auctions:imbalance}{indicative price} wanders by two or three ticks at
first and settles as the paired volume grows.
\end{tutsteps}
\textbf{What to change next.} Add a rule that the \omterm{def:m1:auctions:call}{uncrossing price} must lie
within a collar of the last continuous trade, and extend the call if it does
not --- a \omterm{def:m1:auctions:volint}{volatility interruption} inside the auction. Then add imbalance-only
orders, which may trade only against the surplus.
\end{tutorial}

\begin{omfigure}
\begin{tikzpicture}
\begin{axis}[width=10.5cm,height=4.8cm,
  xlabel={batch of arriving orders},ylabel={indicative price (ticks)},
  tick label style={font=\small},label style={font=\small},
  xmin=1,xmax=40,ymin=996,ymax=1004,grid=major,grid style={gray!25},
  yticklabel style={/pgf/number format/1000 sep={\,}},axis y line*=left,table/col sep=comma]
\addplot[omDef,very thick,const plot] table[x=step,y=price]{figdata/markets-1/13-auctions/indicative.csv};
\end{axis}
\begin{axis}[width=10.5cm,height=4.8cm,
  ylabel={paired volume (thousand shares)},tick label style={font=\small},label style={font=\small},
  xmin=1,xmax=40,ymin=0,ymax=160,axis y line*=right,axis x line=none,table/col sep=comma]
\addplot[omThm,thick,dashed] table[x=step,y=volume_k]{figdata/markets-1/13-auctions/indicative.csv};
\end{axis}
\end{tikzpicture}

\omcaption{A simulated call phase. The \omterm{def:m1:auctions:imbalance}{indicative price} (blue, left axis)
moves freely while little is paired and is pinned once the paired volume
(red dashed, right axis) is large compared with any single order. Data: the
tutorial's simulation.}\label{fig:m1:auctions:indicative}
\end{omfigure}

\section{Build: the auction}\label{bld:m1:auctions:uncross}

\begin{build}
\bfield{Purpose} The exchange simulator of One Quant Book 10 opens and closes
each session with a \omterm{def:m1:auctions:call}{call auction} and interrupts continuous trading with one
when prices move too fast. This is its uncrossing engine.
\bfield{Interface} \texttt{AuctionOrder(order\_id, side, quantity, price,
seq)} with \texttt{price = None} for a market order;
\texttt{uncross(orders, reference)} returning price, volume, signed surplus
and the list of fills; \texttt{demand} and \texttt{supply} for publishing
indicative information.
\bfield{Rules} Integer prices. The four rules of \cref{met:m1:auctions:rules}
in that order; candidate prices are the limits present in the book and the
reference. No cross gives no price and no fills. Bought and sold quantities
in the fills are equal, always.
\bfield{Acceptance tests} \texttt{code/firm/auction/tests/}, \texttt{cpp/} and
\texttt{rust/}: the chapter's example with its partial fill; no cross; buy and
sell pressure; the reference-price tie-break; time priority. The three
implementations must agree.
\bfield{Stretch} Imbalance-only and \omterm{def:m1:auctions:moc}{limit-on-close order} types with entry
cut-offs, and a function that publishes paired quantity, imbalance and
\omterm{def:m1:auctions:imbalance}{indicative price} in the format of \cref{dat:m1:auctions:us}.
\end{build}

\begin{omsources}
\item New York Stock Exchange, \emph{NYSE Closing Process} and \emph{NYSE
Opening and Closing Auctions} fact sheets; \emph{NYSE Closing Auction: Timing
Shifts and Marketability Trends}, 18 November 2025.
\item Nasdaq, \emph{The Nasdaq Opening and Closing Crosses: Frequently Asked
Questions}; US Securities and Exchange Commission, Release 34-84454 (cut-off
times for on-close orders).
\item Deutsche Börse, \emph{Xetra Market Model: Continuous Trading and
Auction}; glossary entry ``Auction principle''.
\item Euronext, \emph{Trading Manual for the Cash Market}.
\item A.~Madhavan, ``Trading mechanisms in securities markets'',
\emph{Journal of Finance} 47 (1992).
\item M.~Pagano and R.~Schwartz, ``A closing call's impact on market
quality at Euronext Paris'', \emph{Journal of Financial Economics} 68
(2003).
\end{omsources}

\section{Exercises}

\begin{exercise}[$\star$]\label{exo:m1:auctions:1}
Buy orders: 400 at 20.10, 300 at 20.05, 500 at 20.00. Sell orders: 200 at
19.95, 600 at 20.05, 300 at 20.10. Tabulate $D$, $S$ and the executable
volume at each limit price and give the \omterm{def:m1:auctions:call}{uncrossing price} and volume.
\end{exercise}

\begin{exercise}[$\star$]\label{exo:m1:auctions:2}
For the auction of the previous exercise give the surplus and its side, and
say which orders are filled, fully or partly.
\end{exercise}

\begin{exercise}[$\star$]\label{exo:m1:auctions:3}
Using \cref{dat:m1:auctions:us}: at 15:53 a trader wants to add a
market-on-close buy order in a stock listed on each exchange. On which can
she, and under what condition on the other?
\end{exercise}

\begin{exercise}[$\star\star$]\label{exo:m1:auctions:4}
A book has a single buy order, 1\,000 shares at 50.20, and a single sell
order, 1\,000 shares at 50.00. The reference price is 50.12. Apply the four
rules. What would the price be with a reference of 50.40?
\end{exercise}

\begin{exercise}[$\star\star$]\label{exo:m1:auctions:5}
Near its \omterm{def:m1:auctions:call}{uncrossing price} an auction has 40\,000 shares of supply and 25\,000
shares of demand per tick. Estimate the effect of an unexpected
market-on-close purchase of 300\,000 shares. The tick is 1 cent and the stock
trades at \$80: express it in basis points.
\end{exercise}

\begin{exercise}[$\star\star$]\label{exo:m1:auctions:6}
An index fund must buy \$50 million of a stock at the close. The published
imbalance suggests the close will be 12 basis points above the current
continuous price; buying now in the continuous market would cost an estimated
8 basis points of impact. What does each choice cost in expectation, and why
will the fund still use the close?
\end{exercise}

\begin{exercise}[$\star\star\star$]\label{exo:m1:auctions:7}
\emph{Coding.} With \texttt{imbalance\_impact}, estimate the slope of the
\omterm{def:m1:auctions:moc}{closing price} with respect to imbalance for books of 400 and of 1\,600
orders (60 books each, imbalance 50\% of the balanced volume). Compare the
slopes \emph{per share of imbalance} and relate them to
\cref{prop:m1:auctions:depth}.
\end{exercise}

\begin{exercise}[$\star\star\star$]\label{exo:m1:auctions:8}
\emph{Find the flaw.} A backtest of a closing-auction strategy assumes that
its \omterm{def:m1:auctions:moc}{limit-on-close orders}, entered at the \omterm{def:m1:auctions:imbalance}{indicative price} shown at 15:55,
are filled in full whenever the \omterm{def:m1:auctions:moc}{closing price} is at or through their limit.
Identify two reasons the fills are overstated, and the information needed to
correct them.
\end{exercise}

\section{Problem: The Index Rebalance Close}

\begin{problem}[{Weekend problem --- a closing auction on rebalance day}]\label{pb:m1:auctions:1}
A stock joins an index at today's close. Ten minutes before the close its
continuous market is $39.99 \times 40.01$ and the last trade, the reference
price, is 40.00; the tick is 1 cent. The closing book holds, besides the
index funds' order, the following.

Buy: 50\,000 at market; 80\,000 at 40.10; 120\,000 at 40.00; 150\,000 at 39.90.
Sell: 60\,000 at market; 100\,000 at 39.95; 140\,000 at 40.05; 200\,000 at
40.15; 250\,000 at 40.30; 300\,000 at 40.50.

\medskip\noindent\textbf{Part I --- Without the index funds.}
\begin{enumerate}
\item Tabulate $D(p)$ and $S(p)$ at each limit price from 39.90 to 40.50.
\item Give the executable volume at each and the \omterm{def:m1:auctions:call}{uncrossing price} and volume.
\item Give the surplus and its side.
\item Which orders are partly filled?
\end{enumerate}

\medskip\noindent\textbf{Part II --- The index funds arrive.}
Index funds add a market-on-close buy order for 500\,000 shares.
\begin{enumerate}[resume]
\item Recompute $D(p)$.
\item Give the new \omterm{def:m1:auctions:call}{uncrossing price} and volume.
\item By how much, in cents and in basis points, did the index order move
the close?
\item What imbalance was published at 15:50 if the index order was already
in, and the uncrossing had been computed at the reference price?
\item How many dollars did the index funds pay above the price of question
2, on their 500\,000 shares?
\end{enumerate}

\medskip\noindent\textbf{Part III --- Liquidity arrives.}
Seeing the imbalance, liquidity providers add sell orders: 150\,000 at 40.10
and 200\,000 at 40.20.
\begin{enumerate}[resume]
\item Recompute $S(p)$ and the \omterm{def:m1:auctions:call}{uncrossing price} and volume.
\item How much did the offsetting orders save the index funds?
\item A provider sold 150\,000 shares at the new close and expects the stock
to revert to 40.05 tomorrow. What is its expected profit?
\item What risk is it taking overnight, and how might it hedge?
\item Estimate the auction's depth $\lambda_D + \lambda_S$ per tick around the
final price from your tables, and check \cref{prop:m1:auctions:depth} on the
move caused by the index order in Part III.
\end{enumerate}

\medskip\noindent\textbf{Part IV --- Judgement.}
\begin{enumerate}[resume]
\item The index funds could have bought during the previous days instead.
What would they have gained and what would they have risked?
\item Who, economically, pays the liquidity providers of question 12?
\item A trader enters a large limit-on-close buy order at 15:57 and cancels
it at 15:59:50 on an exchange that allows it. What was its purpose, and what
is the practice called?
\item Why do exchanges restrict entry and cancellation of closing orders
after the first imbalance publication?
\item State the \emph{named result}: the \omterm{def:m1:auctions:call}{uncrossing price} and the paired
volume of the final auction of Part III.
\item In one sentence: what does an index fund buy with the premium it pays
at the close?
\end{enumerate}
\end{problem}

\section{Interview questions}

\begin{interviewq}[$\star$ \iqroles{trader, developer, researcher}]\label{iq:m1:auctions:1}
How is the price of a closing auction determined?
\end{interviewq}

\begin{interviewq}[$\star$ \iqroles{trader, researcher}]\label{iq:m1:auctions:2}
Why does so much volume trade at the close?
\end{interviewq}

\begin{interviewq}[$\star\star$ \iqroles{developer}]\label{iq:m1:auctions:3}
Implement an uncrossing algorithm that runs in $O(n \log n)$ for $n$ orders.
What are the edge cases?
\end{interviewq}

\begin{interviewq}[$\star\star$ \iqroles{trader, researcher}]\label{iq:m1:auctions:4}
The published closing imbalance in a stock is 800\,000 shares to buy, against
an average auction of 1.5 million. What do you expect to happen in the next
ten minutes, in the auction and in the continuous market?
\end{interviewq}

\begin{interviewq}[$\star\star$ \iqroles{researcher, mle}]\label{iq:m1:auctions:5}
You want to predict the \omterm{def:m1:auctions:moc}{closing price} from imbalance messages. What is your
target, what are your features, and what is the main trap in the
evaluation?
\end{interviewq}

\begin{interviewq}[$\star\star\star$ \iqroles{trader, researcher}]\label{iq:m1:auctions:6}
A venue proposes to replace continuous trading with an auction every 100
milliseconds. Who gains, who loses, and what would you expect to happen to
spreads?
\end{interviewq}
